\documentclass[11pt,a4paper]{article}

\newcommand{\doctitle}{Documento de Casos de Uso}
\newcommand{\docsubtitle}{WARDEN - Wireless Attack Reproduction and Detection Environment}

\usepackage{warden-style}

\begin{document}

% -- Portada -------------------------------------------------------------------
\wardencoverpage{0.1}{Abril 2026}{Borrador}

% -- Historial de Revisiones ---------------------------------------------------
\begin{wardenrevisions}
\revrow{0.1}{Abril 2026}{Santiago Valencia León}{Borrador inicial}
\revrow{0.1}{Abril 2026}{Daniel Yadir Perea Murillo}{Borrador inicial}
\revrow{0.1}{Abril 2026}{Sofia Valdez Londono}{Borrador inicial}
\end{wardenrevisions}

% -- Tabla de Contenidos -------------------------------------------------------
\tableofcontents
\newpage

% ==============================================================================
\section{Introducción}
% ==============================================================================

\subsection{Propósito}
Este documento describe los casos de uso del sistema WARDEN: cada
interacción significativa entre un actor (humano o sistema externo) y el
sistema, junto con sus precondiciones, postcondiciones, flujos principales
y flujos alternos.

Los casos de uso traducen los requerimientos funcionales especificados en
el SRS a escenarios concretos de operación. Cada caso de uso está trazado
a uno o más requerimientos funcionales del SRS.

\subsection{Audiencia}
Este documento está dirigido a los integrantes del equipo de desarrollo
(como guía de qué escenarios deben funcionar) y al profesor evaluador
(como referencia de las interacciones que serán demostradas durante la
entrega).

\subsection{Alcance}
El documento cubre los doce casos de uso identificados a partir del SRS,
agrupados por subsistema (ofensivo, defensivo, sistema completo) y por
actor.

\subsection{Convención de Identificadores}
Cada caso de uso tiene un identificador único con el formato
\texttt{UC-NN}, donde \texttt{NN} es una secuencia de dos dígitos.

\subsection{Estructura de las Fichas}
Cada ficha de caso de uso contiene los siguientes campos estándar:

\begin{datatable}{L{3.2cm} L{12.1cm}}{Campo & Significado}
\drow{ID                        & Identificador único del caso de uso (\texttt{UC-NN}).}
\drow{Nombre                    & Descripción breve y orientada a la acción.}
\drow{Actor Primario            & Actor que inicia el caso de uso.}
\drow{Actores Secundarios       & Otros actores que participan, si aplica.}
\drow{Requerimientos Cubiertos  & IDs de requerimientos del SRS que este caso de uso implementa.}
\drow{Prioridad                 & Importancia del caso de uso (Alta, Media, Baja).}
\drow{Precondiciones            & Condiciones que deben cumplirse antes de iniciar el caso de uso.}
\drow{Postcondiciones           & Estado del sistema al finalizar exitosamente.}
\drow{Flujo Principal           & Secuencia de pasos en el escenario exitoso.}
\drow{Flujos Alternos           & Variaciones y caminos de error.}
\drow{Notas                     & Observaciones técnicas o aclaraciones.}
\end{datatable}

\subsection{Referencias}
\begin{itemize}
  \item Documento WARDEN \emph{Documento de Visión}.
  \item Documento WARDEN \emph{Especificación de Requerimientos de Software}.
  \item Documento WARDEN \emph{Threat Model y Cadena de Ataque}.
\end{itemize}

% ==============================================================================
\section{Actores del Sistema}
% ==============================================================================

WARDEN tiene cuatro actores principales, definidos en el Documento de
Visión y en el SRS:

\begin{datatable}{L{4cm} L{11.3cm}}{Actor & Descripción}
\drow{Operador del Atacante    & Persona que controla la ESP32 y selecciona qué fase ejecutar. Inicia la cadena completa o cada fase individualmente. Interactúa por consola serial o botón físico.}
\drow{Operador del Detector    & Persona que ejecuta el software detector y monitorea las alertas. Interactúa por línea de comandos.}
\drow{Operador de la Víctima   & Persona que opera los dispositivos víctima durante la demostración. Para fines demostrativos, ingresa credenciales ficticias en el portal cautivo.}
\drow{Observador               & Persona que estudia el sistema sin interactuar directamente. Es el rol del profesor evaluador y de los estudiantes espectadores durante la demostración.}
\end{datatable}

% ==============================================================================
\section{Diagrama de Casos de Uso (Vista General)}
% ==============================================================================

El siguiente diagrama representa la relación entre actores y casos de
uso del sistema:

\vspace{0.3cm}
\begin{verbatim}
                 +============================================+
                 |          SISTEMA WARDEN                    |
                 |                                            |
                 |  ----- SUBSISTEMA OFENSIVO -----           |
                 |  ( UC-01: Iniciar cadena automatica )      |
   Operador      |  ( UC-02: Ejecutar Beacon Flooding   )     |
   del         <-+->( UC-03: Ejecutar Deauth            )     |
   Atacante      |  ( UC-04: Ejecutar Evil Twin         )     |
                 |  ( UC-05: Detener ataque en curso    )     |
                 |  ( UC-06: Configurar BSSID objetivo  )     |
                 |                                            |
                 |  ----- PANEL DEL ATACANTE -----            |
                 |  ( UC-13: Reconocer redes WiFi       )     |
   Operador      |  ( UC-14: Reconocer clientes         )     |
   del         <-+->( UC-15: Confirmar etica            )     |
   Atacante      |  ( UC-16: Lanzar ataque desde panel  )     |
                 |  ( UC-17: Monitorear ataque en panel )     |
                 |                                            |
                 |  ----- SUBSISTEMA DEFENSIVO -----          |
                 |  ( UC-07: Iniciar deteccion en vivo  )     |
   Operador      |  ( UC-08: Analizar archivo PCAP      )     |
   del         <-+->( UC-09: Visualizar alertas         )     |
   Detector      |  ( UC-10: Finalizar sesion           )     |
                 |                                            |
                 |  ----- PANEL DEL DEFENSOR -----            |
   Operador      |  ( UC-18: Iniciar Panel del Defensor )     |
   del         <-+->( UC-19: Visualizar amenazas en GUI )     |
   Detector      |  ( UC-20: Reiniciar sesion en panel  )     |
                 |                                            |
                 |  ----- VICTIMA -----                       |
   Operador      |  ( UC-11: Conectarse a Evil Twin e   )     |
   de la       <-+->(        ingresar credenciales      )     |
   Victima       |                                            |
                 |                                            |
                 |  ----- SISTEMA COMPLETO -----              |
   Operador      |                                            |
   del         <-+->( UC-12: Ejecutar demostracion E2E )      |
   Atacante,     |                                            |
   Detector,     |                                            |
   Victima       |                                            |
                 |                                            |
                 +============================================+
\end{verbatim}

% ==============================================================================
\section{Lista de Casos de Uso}
% ==============================================================================

\begin{datatable}{L{1.5cm} L{6.5cm} L{4.5cm} L{1.5cm}}{ID & Nombre & Actor Primario & Prioridad}
\drow{UC-01 & Iniciar cadena ofensiva automática       & Operador del Atacante & Alta}
\drow{UC-02 & Ejecutar Beacon Flooding individualmente & Operador del Atacante & Alta}
\drow{UC-03 & Ejecutar Deauthentication individualmente & Operador del Atacante & Alta}
\drow{UC-04 & Ejecutar Evil Twin individualmente       & Operador del Atacante & Alta}
\drow{UC-05 & Detener ataque en curso                  & Operador del Atacante & Alta}
\drow{UC-06 & Configurar BSSID objetivo                & Operador del Atacante & Alta}
\drow{UC-07 & Iniciar detección en captura en vivo     & Operador del Detector & Alta}
\drow{UC-08 & Analizar archivo PCAP                    & Operador del Detector & Alta}
\drow{UC-09 & Visualizar alertas en tiempo real        & Operador del Detector & Alta}
\drow{UC-10 & Finalizar sesión de detección            & Operador del Detector & Alta}
\drow{UC-11 & Conectarse al Evil Twin e ingresar credenciales ficticias & Operador de la Víctima & Alta}
\drow{UC-12 & Ejecutar demostración end-to-end         & Múltiples actores     & Alta}
\drow{UC-13 & Reconocer redes WiFi cercanas desde el Panel del Atacante & Operador del Atacante & Alta}
\drow{UC-14 & Reconocer clientes asociados a una red e identificar fabricantes & Operador del Atacante & Alta}
\drow{UC-15 & Confirmar autorización ética del objetivo & Operador del Atacante & Alta}
\drow{UC-16 & Lanzar ataque desde el Panel del Atacante & Operador del Atacante & Alta}
\drow{UC-17 & Monitorear ataque en curso desde el Panel del Atacante & Operador del Atacante & Alta}
\drow{UC-18 & Iniciar el Panel del Defensor            & Operador del Detector & Alta}
\drow{UC-19 & Visualizar amenazas detectadas en el Panel del Defensor & Operador del Detector & Alta}
\drow{UC-20 & Reiniciar sesión desde el Panel del Defensor & Operador del Detector & Media}
\end{datatable}

% ==============================================================================
\section{Fichas Detalladas: Subsistema Ofensivo}
% ==============================================================================

\subsection{UC-01: Iniciar cadena ofensiva automática}

\begin{uctable}
\ucrow{ID}{UC-01}
\ucrow{Nombre}{Iniciar cadena ofensiva automática}
\ucrow{Actor Primario}{Operador del Atacante}
\ucrow{Actores Secundarios}{ESP32 (sistema)}
\ucrow{Requerimientos Cubiertos}{RF-OFF-001, RF-OFF-014, RF-OFF-015}
\ucrow{Prioridad}{Alta}
\ucrow{Precondiciones}{
  (1) La ESP32 está alimentada y con el firmware ofensivo cargado. \newline
  (2) El BSSID objetivo está configurado correctamente. \newline
  (3) El router de laboratorio está encendido y emitiendo. \newline
  (4) Al menos un dispositivo víctima está conectado al router de laboratorio.
}
\ucrow{Postcondiciones}{
  Las tres fases de la cadena ofensiva se han ejecutado en secuencia.
  El portal cautivo ha estado activo y, si una víctima ingresó datos,
  estos quedaron registrados en consola serial.
}
\ucrow{Flujo Principal}{
  1. El Operador presiona el botón físico de inicio (o envía el comando
  \texttt{start} por consola serial). \newline
  2. La ESP32 emite por consola serial un mensaje indicando inicio de la
  cadena. \newline
  3. La ESP32 inicia la fase de Beacon Flooding por la duración
  configurada. \newline
  4. Al finalizar Beacon Flooding, la ESP32 emite un mensaje de fin de
  fase 1. \newline
  5. La ESP32 inicia la fase de Deauthentication. \newline
  6. Al finalizar Deauth, la ESP32 emite un mensaje de fin de fase 2. \newline
  7. La ESP32 cambia a modo AP y levanta el Evil Twin con el portal
  cautivo. \newline
  8. La ESP32 espera conexiones de víctimas y registra cualquier credencial
  recibida. \newline
  9. Tras el tiempo configurado o detención manual, la ESP32 finaliza la
  fase 3 y emite un resumen de la cadena.
}
\ucrow{Flujos Alternos}{
  \textbf{A1 - Detención durante la cadena:} En cualquier punto, el
  Operador puede invocar UC-05 (Detener ataque) y la cadena se interrumpe
  emitiendo un resumen parcial. \newline
  \textbf{A2 - BSSID no configurado:} Si en la precondición 2 falla, la
  ESP32 emite un error por consola serial y aborta. \newline
  \textbf{A3 - Ninguna víctima se conecta al Evil Twin:} La fase 3 termina
  sin credenciales registradas. La cadena se considera completa de todas
  formas (la captura es opcional).
}
\ucrow{Notas}{
  La duración de cada fase es configurable según RF-OFF-015. Por defecto:
  Beacon Flood 30 s, Deauth 15 s, Evil Twin 120 s.
}
\end{uctable}

\subsection{UC-02: Ejecutar Beacon Flooding individualmente}

\begin{uctable}
\ucrow{ID}{UC-02}
\ucrow{Nombre}{Ejecutar Beacon Flooding individualmente}
\ucrow{Actor Primario}{Operador del Atacante}
\ucrow{Actores Secundarios}{ESP32 (sistema)}
\ucrow{Requerimientos Cubiertos}{RF-OFF-002, RF-OFF-003, RF-OFF-004}
\ucrow{Prioridad}{Alta}
\ucrow{Precondiciones}{
  (1) La ESP32 está alimentada con el firmware cargado. \newline
  (2) El canal de operación está configurado.
}
\ucrow{Postcondiciones}{
  La ESP32 ha emitido beacons falsos en el canal especificado durante el
  tiempo definido por el Operador.
}
\ucrow{Flujo Principal}{
  1. El Operador envía el comando \texttt{beacon} por consola serial. \newline
  2. La ESP32 emite un mensaje ``Iniciando Beacon Flooding''. \newline
  3. La ESP32 emite al menos 50 beacons únicos por segundo en el canal
  configurado, con SSIDs y BSSIDs falsificados. \newline
  4. La ESP32 mantiene la emisión hasta recibir un comando \texttt{stop} o
  hasta cumplirse la duración configurada. \newline
  5. La ESP32 emite un mensaje ``Beacon Flooding finalizado'' con el conteo
  total de beacons emitidos.
}
\ucrow{Flujos Alternos}{
  \textbf{A1 - Detención manual:} El Operador envía \texttt{stop} antes
  del tiempo configurado y la fase termina inmediatamente.
}
\ucrow{Notas}{
  Los SSIDs falsos se generan a partir de un prefijo configurable o de una
  lista predefinida (RF-OFF-003).
}
\end{uctable}

\subsection{UC-03: Ejecutar Deauthentication individualmente}

\begin{uctable}
\ucrow{ID}{UC-03}
\ucrow{Nombre}{Ejecutar Deauthentication individualmente}
\ucrow{Actor Primario}{Operador del Atacante}
\ucrow{Actores Secundarios}{ESP32 (sistema), Dispositivo víctima}
\ucrow{Requerimientos Cubiertos}{RF-OFF-002, RF-OFF-005, RF-OFF-006}
\ucrow{Prioridad}{Alta}
\ucrow{Precondiciones}{
  (1) La ESP32 está alimentada con el firmware cargado. \newline
  (2) El BSSID del router de laboratorio está configurado en la ESP32. \newline
  (3) El router de laboratorio está encendido. \newline
  (4) Al menos una víctima está asociada al router (ideal para que el
  ataque sea visible).
}
\ucrow{Postcondiciones}{
  La ESP32 ha emitido frames de deauthentication dirigidos al BSSID del
  router. Si había víctimas conectadas, se desconectaron y, en el caso del
  Redmi Note 7, intentaron reconectarse.
}
\ucrow{Flujo Principal}{
  1. El Operador envía el comando \texttt{deauth} por consola serial. \newline
  2. La ESP32 emite ``Iniciando Deauthentication contra BSSID
  \texttt{XX:XX:XX:XX:XX:XX}''. \newline
  3. La ESP32 emite frames de deauth con dirección de origen suplantando
  al BSSID legítimo, dirigidos al broadcast del BSSID o a MACs específicas
  de víctimas. \newline
  4. La emisión continúa hasta el comando \texttt{stop} o el tiempo
  configurado. \newline
  5. La ESP32 emite ``Deauthentication finalizado, frames emitidos: N''.
}
\ucrow{Flujos Alternos}{
  \textbf{A1 - PMF activo en la víctima:} Si el router tiene 802.11w
  (PMF) activado, los deauths son ignorados. La fase termina pero el
  efecto es nulo. Esto se reporta por consola si se detecta. \newline
  \textbf{A2 - Detención manual:} El Operador envía \texttt{stop} antes
  del tiempo configurado.
}
\ucrow{Notas}{
  Para el laboratorio WARDEN, el router se configura sin PMF para
  garantizar que el ataque funcione.
}
\end{uctable}

\subsection{UC-04: Ejecutar Evil Twin individualmente}

\begin{uctable}
\ucrow{ID}{UC-04}
\ucrow{Nombre}{Ejecutar Evil Twin individualmente}
\ucrow{Actor Primario}{Operador del Atacante}
\ucrow{Actores Secundarios}{ESP32 (sistema), Operador de la Víctima}
\ucrow{Requerimientos Cubiertos}{RF-OFF-002, RF-OFF-007, RF-OFF-008, RF-OFF-009, RF-OFF-010, RF-OFF-011, RF-OFF-012}
\ucrow{Prioridad}{Alta}
\ucrow{Precondiciones}{
  (1) La ESP32 está alimentada con el firmware cargado. \newline
  (2) El SSID legítimo a clonar está configurado.
}
\ucrow{Postcondiciones}{
  El Evil Twin estuvo activo durante el tiempo configurado. Si una víctima
  se conectó e ingresó credenciales, estas quedaron registradas con
  marca de tiempo.
}
\ucrow{Flujo Principal}{
  1. El Operador envía el comando \texttt{eviltwin} por consola serial. \newline
  2. La ESP32 emite ``Levantando Evil Twin con SSID
  \texttt{LAB\_WARDEN\_UTP}''. \newline
  3. La ESP32 inicia el AP abierto con el SSID configurado. \newline
  4. La ESP32 inicia el servidor DHCP en la red interna del AP. \newline
  5. La ESP32 inicia el servidor DNS que resuelve cualquier dominio a su
  propia IP. \newline
  6. La ESP32 inicia el servidor web que sirve el portal cautivo. \newline
  7. La ESP32 espera conexiones. Por cada cliente que se conecta, emite
  un mensaje ``Cliente conectado: \texttt{XX:XX:XX:XX:XX:XX} obtuvo IP
  10.0.0.X''. \newline
  8. Si un cliente envía credenciales por el portal, la ESP32 las imprime
  en consola serial con marca de tiempo. \newline
  9. La ESP32 mantiene el Evil Twin hasta el comando \texttt{stop} o el
  tiempo configurado. \newline
  10. Al finalizar, emite ``Evil Twin finalizado, conexiones: N,
  credenciales capturadas: M''.
}
\ucrow{Flujos Alternos}{
  \textbf{A1 - Ningún cliente se conecta:} El Evil Twin termina con N=0
  y M=0. \newline
  \textbf{A2 - Cliente se conecta pero no ingresa credenciales:} El Evil
  Twin termina con N>0 y M=0. \newline
  \textbf{A3 - Detención manual:} El Operador envía \texttt{stop}.
}
\ucrow{Notas}{
  Las credenciales se almacenan únicamente en memoria volátil (RNF-SYS-005).
  La página servida tiene el aviso de demostración educativa visible
  (RF-OFF-011, RNF-SYS-007).
}
\end{uctable}

\subsection{UC-05: Detener ataque en curso}

\begin{uctable}
\ucrow{ID}{UC-05}
\ucrow{Nombre}{Detener ataque en curso}
\ucrow{Actor Primario}{Operador del Atacante}
\ucrow{Actores Secundarios}{ESP32 (sistema)}
\ucrow{Requerimientos Cubiertos}{RF-OFF-013}
\ucrow{Prioridad}{Alta}
\ucrow{Precondiciones}{
  (1) Hay una fase de la cadena ofensiva en ejecución (UC-01, UC-02,
  UC-03 o UC-04 activo).
}
\ucrow{Postcondiciones}{
  La fase en curso ha terminado de forma controlada, liberando recursos
  (modo AP, servidores, emisión de frames).
}
\ucrow{Flujo Principal}{
  1. El Operador envía el comando \texttt{stop} por consola serial o
  presiona el botón físico de detención. \newline
  2. La ESP32 detecta la señal y detiene la emisión de frames o el AP. \newline
  3. La ESP32 cierra los servidores activos (DHCP, DNS, web) si están
  corriendo. \newline
  4. La ESP32 emite ``Detención solicitada, recursos liberados''. \newline
  5. La ESP32 vuelve al estado de espera de comandos.
}
\ucrow{Flujos Alternos}{
  \textbf{A1 - No hay ataque activo:} La ESP32 ignora el comando y emite
  ``Ningún ataque activo''.
}
\ucrow{Notas}{
  Este caso de uso es indispensable como mecanismo de seguridad para
  detener la operación si algo sale mal.
}
\end{uctable}

\subsection{UC-06: Configurar BSSID objetivo}

\begin{uctable}
\ucrow{ID}{UC-06}
\ucrow{Nombre}{Configurar BSSID objetivo}
\ucrow{Actor Primario}{Operador del Atacante}
\ucrow{Actores Secundarios}{ESP32 (sistema)}
\ucrow{Requerimientos Cubiertos}{RF-OFF-006, RNF-SYS-004}
\ucrow{Prioridad}{Alta}
\ucrow{Precondiciones}{
  (1) La ESP32 está alimentada con el firmware cargado.
}
\ucrow{Postcondiciones}{
  El BSSID objetivo queda almacenado en la configuración de la ESP32 y
  será utilizado por las fases subsiguientes.
}
\ucrow{Flujo Principal}{
  1. El Operador envía por consola serial el comando \texttt{set bssid
  XX:XX:XX:XX:XX:XX}. \newline
  2. La ESP32 valida el formato del BSSID. \newline
  3. La ESP32 valida que el BSSID no esté en una lista negra de
  redes externas conocidas (validación de seguridad RNF-SYS-004). \newline
  4. La ESP32 guarda el BSSID en memoria. \newline
  5. La ESP32 confirma ``BSSID objetivo configurado:
  \texttt{XX:XX:XX:XX:XX:XX}''.
}
\ucrow{Flujos Alternos}{
  \textbf{A1 - BSSID con formato inválido:} La ESP32 rechaza el comando y
  emite ``Error: formato de BSSID inválido''. \newline
  \textbf{A2 - BSSID en lista negra:} La ESP32 rechaza el comando y emite
  ``Error: BSSID corresponde a una red externa, configuración rechazada''.
}
\ucrow{Notas}{
  La validación A2 es una restricción ética codificada en el firmware
  para garantizar que el ataque solo se dirija a redes del laboratorio.
}
\end{uctable}

% ==============================================================================
\section{Fichas Detalladas: Panel del Atacante}
% ==============================================================================

\subsection{UC-13: Reconocer redes WiFi cercanas desde el Panel del Atacante}

\begin{uctable}
\ucrow{ID}{UC-13}
\ucrow{Nombre}{Reconocer redes WiFi cercanas desde el Panel del Atacante}
\ucrow{Actor Primario}{Operador del Atacante}
\ucrow{Actores Secundarios}{ESP32 (sistema), Panel del Atacante (frontend)}
\ucrow{Requerimientos Cubiertos}{RF-OFF-016, RF-OFF-017, RF-PAT-001, RF-PAT-002, RF-PAT-003}
\ucrow{Prioridad}{Alta}
\ucrow{Precondiciones}{
  (1) La ESP32 está alimentada con el firmware cargado y en modo control. \newline
  (2) La laptop del Operador está conectada a la red WiFi
  \texttt{WARDEN\_CONTROL} servida por la ESP32. \newline
  (3) El Operador tiene abierto el Panel del Atacante en su navegador.
}
\ucrow{Postcondiciones}{
  El Operador ha visualizado la lista de redes WiFi cercanas y ha
  seleccionado una como objetivo. La selección queda guardada en el
  estado del Panel para los siguientes casos de uso.
}
\ucrow{Flujo Principal}{
  1. El Operador abre el Panel del Atacante en su navegador. \newline
  2. El Panel verifica la conectividad con la ESP32 (\texttt{GET /status}). \newline
  3. El Operador presiona el botón ``Escanear redes''. \newline
  4. El Panel envía \texttt{GET /scan} a la ESP32. \newline
  5. La ESP32 ejecuta un escaneo activo en la banda 2.4 GHz durante
  aproximadamente 10 segundos. \newline
  6. La ESP32 retorna en JSON la lista de redes con SSID, BSSID, canal,
  RSSI y cifrado. \newline
  7. El Panel renderiza la lista en una tabla ordenada por RSSI
  descendente. \newline
  8. El Operador identifica visualmente la red de laboratorio
  \texttt{LAB\_WARDEN\_UTP} (idealmente la de mayor RSSI). \newline
  9. El Operador hace clic en el botón ``Seleccionar'' de esa fila. \newline
  10. El Panel muestra los datos de la red seleccionada en una sección
  destacada y habilita el siguiente paso del flujo (UC-14).
}
\ucrow{Flujos Alternos}{
  \textbf{A1 - ESP32 no responde:} El Panel muestra ``Error de conexión
  con la ESP32'' y sugiere verificar la red \texttt{WARDEN\_CONTROL}. \newline
  \textbf{A2 - Escaneo sin resultados:} El Panel muestra ``No se
  encontraron redes en el rango. Verifique que el router de
  laboratorio esté encendido y dentro del alcance''.
}
\ucrow{Notas}{
  El escaneo activo emite probe requests; esto es comportamiento normal
  de cualquier cliente WiFi y no constituye un ataque.
}
\end{uctable}

\subsection{UC-14: Reconocer clientes asociados a una red e identificar fabricantes}

\begin{uctable}
\ucrow{ID}{UC-14}
\ucrow{Nombre}{Reconocer clientes asociados a una red e identificar fabricantes}
\ucrow{Actor Primario}{Operador del Atacante}
\ucrow{Actores Secundarios}{ESP32 (sistema), Panel del Atacante (frontend)}
\ucrow{Requerimientos Cubiertos}{RF-OFF-018, RF-OFF-019, RF-OFF-024, RF-PAT-004, RF-PAT-005, RF-PAT-006}
\ucrow{Prioridad}{Alta}
\ucrow{Precondiciones}{
  (1) UC-13 completado: hay una red objetivo seleccionada en el Panel. \newline
  (2) Hay al menos un dispositivo víctima conectado a la red objetivo
  generando tráfico (idealmente, el celular Redmi Note 7 o la Dell
  Latitude E6220 navegando o consumiendo datos).
}
\ucrow{Postcondiciones}{
  El Operador ha visualizado la lista de clientes detectados, su
  fabricante por OUI, y ha seleccionado una MAC como víctima. La
  selección queda en el estado del Panel.
}
\ucrow{Flujo Principal}{
  1. El Operador presiona ``Detectar clientes''. \newline
  2. El Panel envía
  \texttt{GET /clients?bssid=XX:XX:XX:XX:XX:XX} a la ESP32. \newline
  3. La ESP32 captura pasivamente frames dirigidos al BSSID indicado
  durante un intervalo configurable (típicamente 30 segundos),
  acumulando las MACs origen distintas. \newline
  4. La ESP32 retorna en JSON la lista de MACs detectadas y el conteo
  de frames observados de cada una. \newline
  5. Para cada MAC de la lista, el Panel ejecuta
  \texttt{GET /oui-lookup?mac=...} contra la ESP32. \newline
  6. La ESP32 consulta su base de datos OUI local y retorna el
  fabricante asociado. \newline
  7. El Panel renderiza una tabla con MAC, fabricante, conteo de frames
  y un botón ``Seleccionar como víctima'' por fila. \newline
  8. El Operador identifica visualmente la víctima esperada (por ejemplo,
  fabricante ``Xiaomi'' para el Redmi Note 7) y la selecciona. \newline
  9. El Panel muestra los datos de la víctima seleccionada y habilita el
  siguiente paso del flujo (UC-15).
}
\ucrow{Flujos Alternos}{
  \textbf{A1 - Sin clientes detectados:} El Panel sugiere prolongar la
  ventana de captura o verificar que las víctimas estén generando
  tráfico. \newline
  \textbf{A2 - OUI desconocido:} Si la base local no contiene el OUI
  consultado, la ESP32 retorna ``Fabricante desconocido'' y el Panel lo
  muestra así. \newline
  \textbf{A3 - El Operador conoce la MAC víctima:} El Operador puede
  ingresar manualmente la MAC en un campo del Panel sin pasar por la
  detección automática.
}
\ucrow{Notas}{
  La identificación por OUI no determina el modelo exacto del dispositivo,
  solo su fabricante. La inferencia humana del Operador (``este es mi
  celular'') es necesaria para confirmar la víctima.
}
\end{uctable}

\subsection{UC-15: Confirmar autorización ética del objetivo}

\begin{uctable}
\ucrow{ID}{UC-15}
\ucrow{Nombre}{Confirmar autorización ética del objetivo}
\ucrow{Actor Primario}{Operador del Atacante}
\ucrow{Actores Secundarios}{Panel del Atacante (frontend)}
\ucrow{Requerimientos Cubiertos}{RF-PAT-007, RNF-SYS-004 (validación ética)}
\ucrow{Prioridad}{Alta}
\ucrow{Precondiciones}{
  (1) UC-13 y UC-14 completados: hay red objetivo y víctima seleccionadas. \newline
  (2) El botón de lanzamiento del ataque permanece deshabilitado hasta
  completar este caso de uso.
}
\ucrow{Postcondiciones}{
  El Operador ha confirmado por escrito que el objetivo es propiedad del
  equipo o cuenta con autorización del propietario. El botón de
  lanzamiento queda habilitado.
}
\ucrow{Flujo Principal}{
  1. El Panel presenta una pantalla intermedia con el resumen del ataque:
  red objetivo, víctima seleccionada, fases que se ejecutarán. \newline
  2. El Panel muestra un texto destacado: ``ADVERTENCIA: estás a punto
  de lanzar un ataque WiFi contra el dispositivo X conectado a la red
  Y. La ejecución de este ataque sobre dispositivos o redes no
  autorizadas constituye un delito según la Ley 1273 de 2009.''. \newline
  3. El Operador debe marcar explícitamente una casilla con el texto:
  ``Confirmo que el objetivo del ataque es propiedad de mi equipo o
  cuento con autorización explícita del propietario''. \newline
  4. Una vez marcada la casilla, el Panel habilita el botón ``Lanzar
  ataque'' (antes deshabilitado en gris). \newline
  5. El Operador puede proceder a UC-16 o cancelar el flujo.
}
\ucrow{Flujos Alternos}{
  \textbf{A1 - Cancelación:} El Operador presiona ``Cancelar'' y el
  Panel vuelve a la pantalla de selección sin habilitar el botón de
  lanzamiento. \newline
  \textbf{A2 - Validador rechaza el BSSID:} Adicionalmente al
  consentimiento del Operador, el firmware ejecuta el Validador Ético
  (UC-06). Si el BSSID está en lista negra, la API rechaza el ataque
  con código HTTP 403 incluso después de la confirmación.
}
\ucrow{Notas}{
  Este caso de uso codifica la responsabilidad ética como parte
  obligatoria del flujo. La doble verificación (consentimiento humano
  en UI + validación técnica en firmware) refuerza el compromiso del
  proyecto con el uso educativo.
}
\end{uctable}

\subsection{UC-16: Lanzar ataque desde el Panel del Atacante}

\begin{uctable}
\ucrow{ID}{UC-16}
\ucrow{Nombre}{Lanzar ataque desde el Panel del Atacante}
\ucrow{Actor Primario}{Operador del Atacante}
\ucrow{Actores Secundarios}{ESP32 (sistema), Panel del Atacante (frontend)}
\ucrow{Requerimientos Cubiertos}{RF-OFF-020, RF-OFF-021, RF-OFF-025, RF-PAT-008, RF-PAT-012}
\ucrow{Prioridad}{Alta}
\ucrow{Precondiciones}{
  (1) UC-13, UC-14 y UC-15 completados. \newline
  (2) El botón de lanzamiento está habilitado.
}
\ucrow{Postcondiciones}{
  La ESP32 ha iniciado la cadena ofensiva contra la víctima
  seleccionada. La transición a UC-17 (monitoreo) es inmediata.
}
\ucrow{Flujo Principal}{
  1. El Operador, opcionalmente, ajusta parámetros del ataque
  (duraciones por fase, modo encadenado vs individual). \newline
  2. El Operador presiona ``Lanzar ataque''. \newline
  3. El Panel envía \texttt{POST /config} con la configuración resultante
  (BSSID, MAC víctima, SSID a clonar, duraciones). \newline
  4. La ESP32 valida la configuración (incluyendo Validador Ético) y
  responde 200 OK o 403 Forbidden según corresponda. \newline
  5. El Panel envía \texttt{POST /attack/start} con el modo de ataque
  seleccionado. \newline
  6. La ESP32 inicia la cadena ofensiva. \newline
  7. El Panel transita automáticamente a la vista de monitoreo (UC-17).
}
\ucrow{Flujos Alternos}{
  \textbf{A1 - Validador ético rechaza:} La ESP32 retorna 403, el
  Panel muestra el mensaje de rechazo y vuelve a la pantalla de
  selección. \newline
  \textbf{A2 - Pérdida de conexión durante el envío:} El Panel muestra
  ``No se recibió confirmación de la ESP32. Verifique manualmente el
  estado por consola serial''.
}
\ucrow{Notas}{
  Este es el último punto donde el Operador puede revisar la
  configuración antes de la ejecución.
}
\end{uctable}

\subsection{UC-17: Monitorear ataque en curso desde el Panel del Atacante}

\begin{uctable}
\ucrow{ID}{UC-17}
\ucrow{Nombre}{Monitorear ataque en curso desde el Panel del Atacante}
\ucrow{Actor Primario}{Operador del Atacante}
\ucrow{Actores Secundarios}{ESP32 (sistema), Panel del Atacante (frontend)}
\ucrow{Requerimientos Cubiertos}{RF-OFF-022, RF-OFF-023, RF-PAT-009, RF-PAT-010, RF-PAT-011, RF-PAT-013}
\ucrow{Prioridad}{Alta}
\ucrow{Precondiciones}{
  (1) UC-16 completado: hay un ataque en curso.
}
\ucrow{Postcondiciones}{
  El Operador ha monitoreado el ataque hasta su finalización (por
  detención manual o por completarse el tiempo configurado). El Panel
  muestra el resumen final.
}
\ucrow{Flujo Principal}{
  1. El Panel muestra la fase actual del ataque, los contadores en vivo
  (beacons emitidos, deauths emitidos, clientes conectados) y un
  cronómetro. \newline
  2. El Panel se actualiza periódicamente (polling cada 1 segundo a
  \texttt{GET /attack/status}, o stream por
  \texttt{GET /events} si está disponible). \newline
  3. Durante la fase 3 (Evil Twin), si llegan credenciales ficticias al
  portal cautivo, el Panel las muestra inmediatamente en una sección
  destacada (consultando \texttt{GET /credentials}). \newline
  4. El botón ``DETENER ATAQUE'' permanece visible y prominente durante
  toda la sesión. \newline
  5. Al completarse la cadena (todas las fases terminadas), el Panel
  muestra la pantalla de resumen post-ataque con duración total,
  contadores finales y credenciales capturadas.
}
\ucrow{Flujos Alternos}{
  \textbf{A1 - Detención manual:} El Operador presiona ``DETENER
  ATAQUE''. El Panel envía \texttt{POST /attack/stop} y la cadena se
  interrumpe. El resumen muestra ``Detenido manualmente''. \newline
  \textbf{A2 - Pérdida temporal de conexión durante Evil Twin:}
  Mientras la ESP32 ejecuta la fase 3, el WiFi
  \texttt{WARDEN\_CONTROL} puede no estar disponible. El Panel muestra
  ``Conexión con ESP32 temporalmente no disponible (fase Evil Twin
  activa)'' y reconecta automáticamente al finalizar la fase. \newline
  \textbf{A3 - Sin credenciales capturadas:} El resumen indica
  ``Ninguna credencial capturada en el portal''.
}
\ucrow{Notas}{
  La pérdida temporal de conectividad durante Evil Twin es esperada y
  documentada. Durante esa fase la ESP32 dedica su radio al AP malicioso.
}
\end{uctable}
% ==============================================================================

\subsection{UC-07: Iniciar detección en captura en vivo}

\begin{uctable}
\ucrow{ID}{UC-07}
\ucrow{Nombre}{Iniciar detección en captura en vivo}
\ucrow{Actor Primario}{Operador del Detector}
\ucrow{Actores Secundarios}{Adaptador Panda PAU0B (sistema), Sistema Operativo Linux}
\ucrow{Requerimientos Cubiertos}{RF-DEF-001, RF-DEF-003, RF-DEF-010, RF-DEF-014}
\ucrow{Prioridad}{Alta}
\ucrow{Precondiciones}{
  (1) El adaptador Panda está conectado y reconocido como \texttt{panda0}. \newline
  (2) El adaptador está en modo monitor (Paso 4 del Manual de Laboratorio
  ejecutado). \newline
  (3) El BSSID y SSID protegidos están configurados como argumentos.
}
\ucrow{Postcondiciones}{
  El detector está capturando frames y aplicando lógica de análisis. La
  consola muestra alertas en tiempo real.
}
\ucrow{Flujo Principal}{
  1. El Operador ejecuta el detector con argumentos: \texttt{python3
  detector.py --iface panda0 --channel 6 --bssid XX:XX:XX:XX:XX:XX
  --ssid LAB\_WARDEN\_UTP}. \newline
  2. El detector valida los argumentos y los muestra por consola. \newline
  3. El detector inicia la captura llamando a Scapy con
  \texttt{sniff(iface=panda0, ...)}. \newline
  4. El detector procesa cada frame entrante aplicando los analizadores
  de cada fase. \newline
  5. Cuando un analizador identifica una firma de ataque, emite una
  alerta por consola (UC-09).
}
\ucrow{Flujos Alternos}{
  \textbf{A1 - Adaptador no en modo monitor:} El detector detecta el
  modo incorrecto y emite ``Error: panda0 no está en modo monitor.
  Ejecute primero los comandos del Manual de Laboratorio Paso 4''. \newline
  \textbf{A2 - Adaptador no encontrado:} El detector emite ``Error:
  interfaz panda0 no encontrada''. \newline
  \textbf{A3 - Argumentos inválidos:} El detector muestra ayuda de uso y
  termina.
}
\ucrow{Notas}{
  La captura requiere privilegios de root o que el ejecutable de Python
  tenga las capabilities apropiadas (\texttt{CAP\_NET\_RAW},
  \texttt{CAP\_NET\_ADMIN}).
}
\end{uctable}

\subsection{UC-08: Analizar archivo PCAP}

\begin{uctable}
\ucrow{ID}{UC-08}
\ucrow{Nombre}{Analizar archivo PCAP}
\ucrow{Actor Primario}{Operador del Detector}
\ucrow{Actores Secundarios}{Sistema de archivos}
\ucrow{Requerimientos Cubiertos}{RF-DEF-002, RF-DEF-003, RNF-DEF-011}
\ucrow{Prioridad}{Alta}
\ucrow{Precondiciones}{
  (1) Existe un archivo PCAP previamente capturado. \newline
  (2) El BSSID y SSID protegidos están configurados.
}
\ucrow{Postcondiciones}{
  El detector ha procesado todos los frames del archivo y emitido las
  alertas correspondientes a las firmas detectadas.
}
\ucrow{Flujo Principal}{
  1. El Operador ejecuta el detector con argumentos: \texttt{python3
  detector.py --pcap captura.pcap --bssid XX:XX:XX:XX:XX:XX --ssid
  LAB\_WARDEN\_UTP}. \newline
  2. El detector verifica que el archivo existe y es legible. \newline
  3. El detector procesa cada frame del archivo aplicando los
  analizadores. \newline
  4. Por cada firma detectada, emite una alerta. \newline
  5. Al finalizar el archivo, emite el resumen de sesión (UC-10).
}
\ucrow{Flujos Alternos}{
  \textbf{A1 - Archivo no existe:} El detector emite ``Error: archivo
  captura.pcap no encontrado'' y termina. \newline
  \textbf{A2 - Archivo corrupto:} El detector emite ``Error: archivo PCAP
  corrupto o formato no soportado''.
}
\ucrow{Notas}{
  Este caso de uso permite que los integrantes del equipo desarrollen y
  prueben el detector sin necesidad de tener el adaptador Panda
  conectado, usando capturas de referencia generadas previamente
  (RF-SYS-002).
}
\end{uctable}

\subsection{UC-09: Visualizar alertas en tiempo real}

\begin{uctable}
\ucrow{ID}{UC-09}
\ucrow{Nombre}{Visualizar alertas en tiempo real}
\ucrow{Actor Primario}{Operador del Detector}
\ucrow{Actores Secundarios}{Detector (sistema)}
\ucrow{Requerimientos Cubiertos}{RF-DEF-008, RF-DEF-009, RF-DEF-007, RNF-DEF-002, RNF-DEF-004}
\ucrow{Prioridad}{Alta}
\ucrow{Precondiciones}{
  (1) El detector está en ejecución (UC-07 o UC-08 activo).
}
\ucrow{Postcondiciones}{
  El Operador ha visto en consola las alertas correspondientes a los
  ataques detectados.
}
\ucrow{Flujo Principal}{
  1. El detector identifica una firma de ataque en el flujo de captura. \newline
  2. El detector construye un mensaje de alerta con marca de tiempo ISO
  8601, severidad, tipo de ataque y metadatos relevantes. \newline
  3. El detector emite el mensaje por la salida estándar. \newline
  4. El Operador lee el mensaje en su consola. \newline
  5. Si tres alertas (Beacon Flood, Deauth, Evil Twin) ocurren en una
  ventana temporal corta, el detector emite una alerta agregada de
  ``Cadena ofensiva detectada''.
}
\ucrow{Flujos Alternos}{
  \textbf{A1 - Falso positivo:} El Operador identifica que la alerta no
  corresponde a un ataque real (típicamente un AP vecino con SSID
  parecido). El detector permite agregar a lista blanca (RF-DEF-015) en
  ejecuciones futuras.
}
\ucrow{Notas}{
  El formato de alerta es legible directamente por humanos. Ejemplo: \newline
  \texttt{[2026-04-25T14:23:01] [ALERT] EVIL\_TWIN BSSID
  AA:BB:CC:DD:EE:FF announces SSID LAB\_WARDEN\_UTP, but protected BSSID
  is XX:XX:XX:XX:XX:XX.}
}
\end{uctable}

\subsection{UC-10: Finalizar sesión de detección}

\begin{uctable}
\ucrow{ID}{UC-10}
\ucrow{Nombre}{Finalizar sesión de detección}
\ucrow{Actor Primario}{Operador del Detector}
\ucrow{Actores Secundarios}{Detector (sistema)}
\ucrow{Requerimientos Cubiertos}{RF-DEF-012, RF-DEF-013}
\ucrow{Prioridad}{Alta}
\ucrow{Precondiciones}{
  (1) El detector está en ejecución.
}
\ucrow{Postcondiciones}{
  El detector ha emitido un resumen de sesión y ha liberado los recursos
  (interfaz de captura, búferes, etc.).
}
\ucrow{Flujo Principal}{
  1. El Operador presiona Ctrl+C en la consola. \newline
  2. El detector intercepta la señal SIGINT. \newline
  3. El detector detiene la captura de Scapy. \newline
  4. El detector calcula el resumen: duración total, conteo de alertas
  por tipo (Beacon Flood, Deauth, Evil Twin, Cadena), conteo de frames
  procesados. \newline
  5. El detector imprime el resumen por consola. \newline
  6. El detector termina con código de salida 0.
}
\ucrow{Flujos Alternos}{
  \textbf{A1 - Finalización por fin de PCAP:} En modo UC-08, la sesión
  termina automáticamente al procesar el último frame del archivo. El
  resumen se emite igual.
}
\ucrow{Notas}{
  El manejo correcto de SIGINT es obligatorio para no dejar la interfaz
  en estado inconsistente.
}
\end{uctable}

% ==============================================================================
\section{Fichas Detalladas: Panel del Defensor}
% ==============================================================================

\subsection{UC-18: Iniciar el Panel del Defensor}

\begin{uctable}
\ucrow{ID}{UC-18}
\ucrow{Nombre}{Iniciar el Panel del Defensor}
\ucrow{Actor Primario}{Operador del Detector}
\ucrow{Actores Secundarios}{Servidor FastAPI (sistema), Adaptador Panda PAU0B (hardware)}
\ucrow{Requerimientos Cubiertos}{RF-PDE-001, RF-PDE-002, RF-PDE-004, RF-PDE-008}
\ucrow{Prioridad}{Alta}
\ucrow{Precondiciones}{
  (1) El adaptador Panda está conectado y reconocido como
  \texttt{panda0}. \newline
  (2) Las dependencias Python (FastAPI, Uvicorn, Scapy) están
  instaladas. \newline
  (3) El BSSID y SSID protegidos están disponibles para configurar.
}
\ucrow{Postcondiciones}{
  El servidor FastAPI está corriendo en \texttt{localhost:8000}, el
  detector está integrado y listo para iniciar captura, el navegador
  muestra el Panel del Defensor.
}
\ucrow{Flujo Principal}{
  1. El Operador ejecuta \texttt{python3 -m warden.defender.web} en su
  terminal. \newline
  2. El servidor FastAPI inicia en el puerto 8000 e imprime la URL de
  acceso. \newline
  3. El Operador abre \texttt{http://localhost:8000} en su
  navegador. \newline
  4. El Panel carga la pantalla principal mostrando el estado del
  detector (``Detenido''), un formulario de configuración (interfaz,
  canal, BSSID protegido, SSID protegido) y un botón ``Iniciar
  detección''. \newline
  5. El Operador ingresa los parámetros de configuración (o acepta los
  valores por defecto cargados del archivo de configuración). \newline
  6. El Operador presiona ``Iniciar detección''. \newline
  7. El servidor verifica que la interfaz esté en modo monitor; si no,
  muestra error sugiriendo ejecutar el procedimiento del Manual de
  Laboratorio. \newline
  8. El servidor inicia el detector en un hilo de fondo y abre la
  conexión WebSocket con el frontend. \newline
  9. El Panel cambia el indicador de estado a ``Corriendo'' y muestra el
  indicador visual de amenaza en verde. \newline
  10. El Panel transita a la vista de monitoreo (UC-19).
}
\ucrow{Flujos Alternos}{
  \textbf{A1 - Puerto 8000 ocupado:} El servidor falla al arrancar; el
  Operador puede pasar argumento \texttt{--port} para usar otro
  puerto. \newline
  \textbf{A2 - Adaptador no en modo monitor:} El servidor responde con
  un mensaje de error en el Panel sugiriendo activar modo monitor
  primero. \newline
  \textbf{A3 - Configuración inválida (BSSID mal formado):} El Panel
  muestra error de validación y conserva los valores ingresados para
  corrección.
}
\ucrow{Notas}{
  La integración del detector dentro del servidor FastAPI permite que
  ambos compartan estado en memoria sin necesidad de comunicación
  entre procesos. El detector original sigue ejecutándose como un
  componente integrado.
}
\end{uctable}

\subsection{UC-19: Visualizar amenazas detectadas en el Panel del Defensor}

\begin{uctable}
\ucrow{ID}{UC-19}
\ucrow{Nombre}{Visualizar amenazas detectadas en el Panel del Defensor}
\ucrow{Actor Primario}{Operador del Detector}
\ucrow{Actores Secundarios}{Detector (componente interno), Servidor FastAPI, WebSocket}
\ucrow{Requerimientos Cubiertos}{RF-PDE-003, RF-PDE-005, RF-PDE-006, RF-PDE-007, RF-PDE-010, RNF-PDE-001, RNF-PDE-002}
\ucrow{Prioridad}{Alta}
\ucrow{Precondiciones}{
  (1) UC-18 completado: el Panel está activo y el detector corriendo.
}
\ucrow{Postcondiciones}{
  El Operador ha visualizado las alertas en tiempo real y el indicador
  visual ha reflejado el nivel de amenaza correspondiente.
}
\ucrow{Flujo Principal}{
  1. El detector procesa frames y, al identificar una firma, construye
  una alerta. \newline
  2. El detector entrega la alerta al servidor FastAPI mediante el
  patrón observador o cola interna. \newline
  3. El servidor envía la alerta al frontend a través de la conexión
  WebSocket establecida. \newline
  4. El frontend recibe la alerta y la inserta al tope de la lista de
  alertas, con animación de aparición. \newline
  5. El frontend incrementa el contador agregado correspondiente al
  tipo de alerta. \newline
  6. El frontend evalúa el indicador visual: si la alerta es de tipo
  \texttt{CADENA\_OFENSIVA}, el indicador se pone en rojo; si es la
  primera alerta de cualquier otro tipo, el indicador transita de verde
  a amarillo. \newline
  7. El Operador observa visualmente el cambio de estado y la nueva
  entrada en la lista. \newline
  8. El Operador puede expandir cualquier alerta para ver sus
  metadatos completos (BSSID, MAC origen, contadores).
}
\ucrow{Flujos Alternos}{
  \textbf{A1 - Pérdida de conexión WebSocket:} El frontend detecta la
  desconexión, muestra ``Conexión perdida con el detector'' y reintenta
  reconectarse cada 3 segundos. \newline
  \textbf{A2 - Múltiples alertas simultáneas:} El frontend acumula las
  alertas en el orden recibido y respeta el orden cronológico. \newline
  \textbf{A3 - Alerta duplicada:} El detector ya implementa cooldown
  por tipo de alerta, así que el frontend confía en la deduplicación
  del backend.
}
\ucrow{Notas}{
  La salida por consola estándar (RF-DEF-008) sigue activa en paralelo:
  el Operador puede ver las mismas alertas tanto en el navegador como
  en la terminal donde ejecuta el servidor.
}
\end{uctable}

\subsection{UC-20: Reiniciar sesión desde el Panel del Defensor}

\begin{uctable}
\ucrow{ID}{UC-20}
\ucrow{Nombre}{Reiniciar sesión desde el Panel del Defensor}
\ucrow{Actor Primario}{Operador del Detector}
\ucrow{Actores Secundarios}{Servidor FastAPI, Detector}
\ucrow{Requerimientos Cubiertos}{RF-PDE-009}
\ucrow{Prioridad}{Media}
\ucrow{Precondiciones}{
  (1) El Panel está activo (con o sin alertas previas).
}
\ucrow{Postcondiciones}{
  La lista de alertas y los contadores agregados quedan limpios. El
  detector sigue corriendo (no se detiene). El indicador visual vuelve
  a verde.
}
\ucrow{Flujo Principal}{
  1. El Operador presiona el botón ``Reiniciar sesión'' en el Panel. \newline
  2. El Panel muestra una confirmación: ``¿Limpiar todas las alertas y
  contadores actuales? El detector seguirá corriendo''. \newline
  3. El Operador confirma. \newline
  4. El frontend envía
  \texttt{POST /api/session/reset} al servidor. \newline
  5. El servidor limpia los contadores y la cola de alertas en memoria. \newline
  6. El servidor emite por WebSocket un evento de reset. \newline
  7. El frontend limpia la lista visual de alertas, resetea los
  contadores a cero y vuelve el indicador a verde.
}
\ucrow{Flujos Alternos}{
  \textbf{A1 - Cancelación:} El Operador cancela el diálogo y la sesión
  continúa intacta.
}
\ucrow{Notas}{
  Este caso de uso es útil entre demostraciones consecutivas para
  empezar con estado limpio sin reiniciar todo el servicio.
}
\end{uctable}

% ==============================================================================
\section{Fichas Detalladas: Víctima}
% ==============================================================================

\subsection{UC-11: Conectarse al Evil Twin e ingresar credenciales ficticias}

\begin{uctable}
\ucrow{ID}{UC-11}
\ucrow{Nombre}{Conectarse al Evil Twin e ingresar credenciales ficticias}
\ucrow{Actor Primario}{Operador de la Víctima}
\ucrow{Actores Secundarios}{Dispositivo víctima (Redmi Note 7 o Dell Latitude E6220), ESP32 (sistema)}
\ucrow{Requerimientos Cubiertos}{RF-OFF-007, RF-OFF-008, RF-OFF-010, RF-OFF-012, RNF-SYS-005, RNF-SYS-007}
\ucrow{Prioridad}{Alta}
\ucrow{Precondiciones}{
  (1) El Evil Twin está activo (UC-04 en curso). \newline
  (2) El dispositivo víctima está dentro del rango de cobertura de la
  ESP32. \newline
  (3) El dispositivo víctima está desconectado del router legítimo (típicamente
  por efecto previo de UC-03).
}
\ucrow{Postcondiciones}{
  El dispositivo víctima ha completado el ciclo: conexión al Evil Twin,
  visualización del portal cautivo, ingreso de credenciales ficticias y
  envío del formulario. La ESP32 ha registrado las credenciales.
}
\ucrow{Flujo Principal}{
  1. El Operador de la Víctima abre la lista de redes WiFi en el
  dispositivo. \newline
  2. El Operador identifica la red \texttt{LAB\_WARDEN\_UTP} (sin
  candado, abierta) y selecciona conectarse. \newline
  3. El dispositivo obtiene IP por DHCP del Evil Twin. \newline
  4. El sistema operativo del dispositivo (Android en el Redmi Note 7)
  detecta la red abierta con portal cautivo y abre el portal
  automáticamente, o el Operador navega a cualquier URL para
  desencadenarlo. \newline
  5. El dispositivo recibe la página HTML del portal cautivo. \newline
  6. El Operador verifica el aviso visible de ``DEMOSTRACION EDUCATIVA''. \newline
  7. El Operador ingresa las credenciales ficticias estándar definidas en
  el Threat Model: \newline
  \quad - Usuario: \texttt{demo.user@warden.test} \newline
  \quad - Contraseña: \texttt{demo-password-12345} \newline
  8. El Operador pulsa el botón de envío. \newline
  9. El navegador envía un POST con los datos al servidor de la ESP32. \newline
  10. La ESP32 registra las credenciales en consola serial con marca de
  tiempo y dirección IP del cliente.
}
\ucrow{Flujos Alternos}{
  \textbf{A1 - El portal no se abre automáticamente:} El Operador abre
  un navegador y va a una URL cualquiera. El DNS de la ESP32 redirige a
  la página del portal. \newline
  \textbf{A2 - El dispositivo no se conecta al Evil Twin:} Posibles
  causas: el dispositivo prefiere reconectarse al router legítimo (si
  está disponible), señal débil, o limitaciones del Android sobre redes
  abiertas. El Operador puede ``olvidar'' la red legítima en la
  configuración del dispositivo para forzar la conexión al Evil Twin. \newline
  \textbf{A3 - Cancelación voluntaria:} El Operador puede cerrar el
  portal sin enviar credenciales. La fase termina sin captura.
}
\ucrow{Notas}{
  Bajo ninguna circunstancia se ingresan credenciales reales en este
  caso de uso. Las credenciales ficticias estándar están definidas en el
  Threat Model como activos demostrativos del proyecto.
}
\end{uctable}

% ==============================================================================
\section{Fichas Detalladas: Sistema Completo}
% ==============================================================================

\subsection{UC-12: Ejecutar demostración end-to-end}

\begin{uctable}
\ucrow{ID}{UC-12}
\ucrow{Nombre}{Ejecutar demostración end-to-end}
\ucrow{Actor Primario}{Operador del Atacante (coordinador)}
\ucrow{Actores Secundarios}{Operador del Detector, Operador de la Víctima, Observador (profesor evaluador)}
\ucrow{Requerimientos Cubiertos}{RF-SYS-001, todos los RF cubiertos por UC-01 a UC-11}
\ucrow{Prioridad}{Alta}
\ucrow{Precondiciones}{
  (1) Todos los componentes físicos del laboratorio están desplegados:
  router de laboratorio, dispositivos víctima, ESP32 con firmware,
  laptop detector con Panda. \newline
  (2) Las víctimas están conectadas al router legítimo. \newline
  (3) El Panda está en modo monitor en el canal del router. \newline
  (4) El detector está iniciado y esperando.
}
\ucrow{Postcondiciones}{
  La cadena ofensiva completa fue ejecutada y el detector identificó
  cada fase, así como la cadena agregada. Las credenciales ficticias de
  demostración fueron capturadas. El Observador presenció la sesión.
}
\ucrow{Flujo Principal}{
  1. El Operador del Detector ejecuta UC-07 (iniciar detección). El
  detector queda escuchando. \newline
  2. El Observador verifica que el detector está activo y que la línea
  base de tráfico está limpia (sin alertas espurias). \newline
  3. El Operador del Atacante ejecuta UC-01 (iniciar cadena automática). \newline
  4. La ESP32 inicia la fase 1 (Beacon Flooding). \newline
  5. El detector identifica la firma de Beacon Flooding y emite la
  alerta correspondiente (UC-09). \newline
  6. La ESP32 transita a la fase 2 (Deauthentication). \newline
  7. El detector identifica el deauth y emite la alerta correspondiente. \newline
  8. La víctima Redmi Note 7 se desconecta del router legítimo. \newline
  9. La ESP32 transita a la fase 3 (Evil Twin). \newline
  10. El detector identifica el Evil Twin y emite la alerta
  correspondiente. \newline
  11. El Operador de la Víctima ejecuta UC-11 (conectarse e ingresar
  credenciales ficticias). \newline
  12. La ESP32 captura las credenciales y las muestra por consola
  serial. \newline
  13. El detector emite la alerta agregada de ``Cadena ofensiva
  detectada'' tras correlacionar las tres fases. \newline
  14. El Operador del Atacante deja terminar la cadena o invoca UC-05. \newline
  15. El Operador del Detector ejecuta UC-10 (finalizar sesión y ver
  resumen). \newline
  16. El Observador revisa los logs de ambos lados y verifica que la
  cadena ofensiva y la detección coinciden temporalmente.
}
\ucrow{Flujos Alternos}{
  \textbf{A1 - Falla parcial de detección:} Si el detector pierde alguna
  fase, las demás se reportan normalmente. La cadena agregada solo se
  emite si las tres fases fueron detectadas. \newline
  \textbf{A2 - Falla del Evil Twin:} Si la víctima no se conecta al Evil
  Twin (flujo alterno A2 de UC-11), la fase 3 termina sin captura. La
  cadena se considera ejecutada igualmente; solo el resultado final
  cambia. \newline
  \textbf{A3 - Detención de emergencia:} Si algo sale mal (por ejemplo,
  un dispositivo ajeno se ve afectado), cualquier Operador puede
  invocar UC-05 y la demostración se aborta. Esto se documenta en la
  bitácora de la sesión (RF-SYS-005).
}
\ucrow{Notas}{
  Este es el caso de uso más importante del proyecto: corresponde a la
  demostración que los evaluadores presenciarán durante la entrega
  académica. Su éxito determina la aceptación del sistema según los
  criterios definidos en el SRS.
}
\end{uctable}

% ==============================================================================
\section{Trazabilidad: Casos de Uso a Requerimientos}
% ==============================================================================

\begin{datatable}{L{1.5cm} L{14cm}}{Caso de Uso & Requerimientos del SRS Cubiertos}
\drow{UC-01 & RF-OFF-001, RF-OFF-014, RF-OFF-015}
\drow{UC-02 & RF-OFF-002, RF-OFF-003, RF-OFF-004}
\drow{UC-03 & RF-OFF-002, RF-OFF-005, RF-OFF-006}
\drow{UC-04 & RF-OFF-002, RF-OFF-007, RF-OFF-008, RF-OFF-009, RF-OFF-010, RF-OFF-011, RF-OFF-012}
\drow{UC-05 & RF-OFF-013}
\drow{UC-06 & RF-OFF-006, RNF-SYS-004}
\drow{UC-07 & RF-DEF-001, RF-DEF-003, RF-DEF-010, RF-DEF-014}
\drow{UC-08 & RF-DEF-002, RF-DEF-003, RNF-DEF-011}
\drow{UC-09 & RF-DEF-007, RF-DEF-008, RF-DEF-009, RNF-DEF-002, RNF-DEF-004}
\drow{UC-10 & RF-DEF-012, RF-DEF-013}
\drow{UC-11 & RF-OFF-007, RF-OFF-008, RF-OFF-010, RF-OFF-012, RNF-SYS-005, RNF-SYS-007}
\drow{UC-12 & RF-SYS-001 (y todos los RF cubiertos por UC-01 a UC-11)}
\drow{UC-13 & RF-OFF-016, RF-OFF-017, RF-PAT-001, RF-PAT-002, RF-PAT-003}
\drow{UC-14 & RF-OFF-018, RF-OFF-019, RF-OFF-024, RF-PAT-004, RF-PAT-005, RF-PAT-006}
\drow{UC-15 & RF-PAT-007, RNF-SYS-004}
\drow{UC-16 & RF-OFF-020, RF-OFF-021, RF-OFF-025, RF-PAT-008, RF-PAT-012}
\drow{UC-17 & RF-OFF-022, RF-OFF-023, RF-PAT-009, RF-PAT-010, RF-PAT-011, RF-PAT-013}
\drow{UC-18 & RF-PDE-001, RF-PDE-002, RF-PDE-004, RF-PDE-008}
\drow{UC-19 & RF-PDE-003, RF-PDE-005, RF-PDE-006, RF-PDE-007, RF-PDE-010, RNF-PDE-001, RNF-PDE-002}
\drow{UC-20 & RF-PDE-009}
\end{datatable}

\end{document}
